\sgaStableTarget{sga6.e13.expose}\section*{\underline{EXPOS\'E XIII}\\[3pt] \underline{\underline{FINITENESS THEOREMS FOR THE PICARD FUNCTOR}}}
\begin{center}
\underline{by S. KLEIMAN}
\end{center}

\sgaStableTarget{sga6.e13.sec.1}\subsection*{1. The (b)-sheaves}
Let @@M0@@ be an algebraically closed field and @@M1@@ a projective @@M2@@-scheme, endowed with an ample invertible module @@M3@@.

\sgaStableTarget{sga6.e13.def.1-1}\textbf{Definition 1.1 (Mumford).} \emph{Let @@M4@@. A coherent sheaf @@M5@@ on @@M6@@ is said to be \emph{@@M7@@-regular} if, at every point of the support of @@M8@@, @@M9@@ is generated by its global sections and if, for every @@M10@@, we have}
@@M11@@

\sgaStableTarget{sga6.e13.lem.1-2}\textbf{Lemma 1.2.} \emph{Let @@M12@@ be an exact sequence of coherent sheaves on @@M13@@}\SGAAuthorFootnote{*}{Here, as below in analogous situations, it is understood that @@M14@@ is defined by tensoring with a given section of @@M15@@.}\emph{. If @@M16@@ is @@M17@@-regular, then so is @@M18@@.}

Indeed, we have the corresponding exact sequence:
@@M19@@

\sgaStableTarget{sga6.e13.prop.1-3}\textbf{Proposition 1.3.} \emph{If a sheaf @@M20@@ on @@M21@@ is @@M22@@-regular, then, for every @@M23@@, we have}~:
\begin{enumerate}[label=(\roman*)]
\item \emph{@@M24@@ is @@M25@@-regular.}
\item \emph{@@M26@@ is surjective.}
\item \emph{@@M27@@ is generated by @@M28@@.}
\end{enumerate}

